\begin{theorem}[原始覆盖下的双时刻可观测二歧]\label{thm:pom-elliptic-two-time-observability-dichotomy}
设 \(k\) 为特征 \(0\) 的代数闭域，\(E/k\) 为椭圆曲线，\(\varphi:E\to E\) 为次数 \(\deg\varphi>1\) 的等ogeny（例如乘法同态 \([n]\)）。
设 \(f:E\to \PP^{1}\) 为次数 \(d\ge 2\) 的有限态射，并令函数域 \(K:=k(E)\)、\(F:=k(f)\subset K\)。
假设扩张 \(K/F\) 为原始扩张，即不存在真中间域 \(F\subsetneq F'\subsetneq K\)。
则恰有下列二者之一成立：
\begin{enumerate}
  \item 存在有理映射 \(R\in k(z)\) 使
  \[
  f\circ\varphi = R\circ f,
  \]
  等价地 \(f\circ\varphi\in F\)。
  此时由观测递推 \(y_{t+1}=R(y_t)\) 可在 \(\PP^1\) 上闭合表示。
  \item 令
  \[
  \Phi:E\to \PP^{1}\times\PP^{1},\qquad \Phi(P):=\bigl(f(P),\,f(\varphi(P))\bigr),
  \]
  则 \(\Phi\) 在其像上为双有理同构；等价地，
  \[
  K = F\bigl(f\circ\varphi\bigr).
  \]
  因而在去掉一闭代数子集后，隐藏态 \(P\) 可由相邻两次观测 \((y_t,y_{t+1})\) 唯一确定。
\end{enumerate}
\leanverified{paper\_pom\_elliptic\_two\_time\_observability\_dichotomy}
\end{theorem}
\begin{proof}
设 \(F':=F(f\circ\varphi)\)。则有域塔 \(F\subseteq F'\subseteq K\)。
由原始性，必有 \(F'=F\) 或 \(F'=K\)。
若 \(F'=F\)，则 \(f\circ\varphi\in F=k(f)\)，故存在 \(R\in k(z)\) 使 \(f\circ\varphi=R\circ f\)。
若 \(F'=K\)，则 \(K=F(f\circ\varphi)\)，这等价于 \(\Phi\) 在其像上诱导函数域同构，从而 \(\Phi\) 为双有理。
二者互斥且覆盖全部情形，结论成立。
\end{proof}

\begin{theorem}[纤维迹矩的有理性与均方最优预测的闭式表达]\label{thm:pom-fiber-trace-moments-rationality}
设 \(C/k\) 为光滑射影曲线，\(f:C\to\PP^1\) 为次数 \(d\) 的有限态射。
令 \(K:=k(C)\)、\(F:=k(y)\)（其中 \(y\) 为 \(\PP^1\) 的坐标，亦即 \(F=k(f)\)）。
任取 \(h\in K\)。
对每个整数 \(m\ge 1\)，在 \(f\) 的分支值之外定义纤维幂和
\[
S_m(y):=\sum_{P\in f^{-1}(y)} h(P)^m\qquad(\text{按重数计}).
\]
则 \(S_m\) 延拓为 \(F\) 中有理函数，并且满足
\[
S_m=\Tr_{K/F}(h^m)\in F.
\]
因而均值与方差函数
\[
\mu(y):=\frac{1}{d}S_1(y),\qquad
v(y):=\frac{1}{d}S_2(y)-\left(\frac{1}{d}S_1(y)\right)^{\!2}
\]
均为 \(y\) 的有理函数。
\par
进一步，若给定递归 \(\varphi:C\to C\) 并取 \(h=f\circ\varphi\)，则在“只观测到底点 \(y_t=f(P_t)\) 且对纤维 \(f^{-1}(y_t)\) 采用均匀后验”的模型下，
\(\mu(y_t)\) 给出对下一步观测 \(y_{t+1}=f(\varphi(P_t))\) 的最优 \(L^2\) 预测，
而 \(v(y_t)\) 为该模型下不可约的均方误差，从而提供误差控制的闭式表达。
\leanverified{paper\_pom\_fiber\_trace\_moments\_rationality}
\end{theorem}
\begin{proof}
由于 \(K/F\) 为有限可分扩张，对任意 \(m\ge 1\) 有 \(\Tr_{K/F}(h^m)\in F\)。
另一方面，在避开分支值的 \(y\) 上，几何纤维 \(f^{-1}(y)\)（按重数）与 \(K\) 在 \(\overline F\) 上的 \(F\)-嵌入所对应的共轭点一致，
迹可解释为这些共轭取值之和，故与 \(S_m(y)\) 在一般点一致，从而二者给出同一 \(F\)-有理函数。
\(\mu,v\) 的有理性由代数运算直接推出。
最后一段为有限离散均匀条件分布下的标准计算：最优 \(L^2\) 预测为条件期望，均方误差为条件方差；此处条件期望与条件二阶矩分别由纤维平均与纤维二阶矩给出，即 \(\mu\) 与 \(v\)。
\end{proof}

\begin{proposition}[零方差判别与模闭合的等价]\label{prop:pom-zero-variance-iff-closure}
在定理 \ref{thm:pom-fiber-trace-moments-rationality} 的记号下，设 \(v(y)\) 由 \(h\in K\) 定义。
则
\[
v\equiv 0\ \text{于}\ F
\quad\Longleftrightarrow\quad
h\in F.
\]
特别地，对 \(h=f\circ\varphi\) 有
\[
v\equiv 0
\quad\Longleftrightarrow\quad
\exists\,R\in k(z)\ \text{使得}\ f\circ\varphi=R\circ f,
\]
亦即观测递归在 \(\PP^1\) 上闭合且无需外置记忆。
\leanverified{paper\_pom\_zero\_variance\_iff\_closure}
\end{proposition}
\begin{proof}
若 \(h\in F\)，则 \(h\) 在同一纤维上为常值函数，方差恒为零。
反之若 \(v\equiv 0\)，则对一般 \(y\) 的纤维，\(\{h(P)\}_{P\in f^{-1}(y)}\) 为常值多重集，
这等价于 \(h\) 被所有 \(F\)-嵌入固定，从而 \(h\in F\)。
代入 \(h=f\circ\varphi\) 即得后半结论。
\end{proof}

\begin{proposition}[椭圆环面的圆维与最小双通道可观测性]\label{prop:pom-elliptic-circle-dimension-minimality}
\leanverified{paper\_pom\_elliptic\_circle\_dimension\_minimality}
将 \(E(\CC)\) 视为紧连通阿贝尔李群 \(E(\CC)\cong \RR^2/\Lambda\cong \TT^2\)，其圆维为 \(2\)。
则任一非常值连续观测 \(g:E(\CC)\to S^1\) 不可能为单射。
并且在定理 \ref{thm:pom-elliptic-two-time-observability-dichotomy} 的第二种情形下，
二时刻观测映射 \(P\mapsto\bigl(f(P),f(\varphi(P))\bigr)\) 在一般点上可逆，
从而两次观测所提供的两条圆通道在一般点上达到恢复态所需的最小圆因子数 \(2\)。
\end{proposition}
\begin{proof}
若存在连续单射 \(E(\CC)\hookrightarrow S^1\)，则其像为 \(S^1\) 的紧子集并与 \(E(\CC)\) 同胚；
但 \(E(\CC)\) 为二维紧流形而 \(S^1\) 为一维流形，矛盾。
第二段由定理 \ref{thm:pom-elliptic-two-time-observability-dichotomy} 的双有理可逆性给出一般点上的二时刻重构；
而一时刻圆观测不可能单射，故两条圆通道在维数意义下达到最小。
\end{proof}

\begin{proposition}[双时刻重构的局部误差放大界]\label{prop:pom-two-time-reconstruction-condition-number}
处于定理 \ref{thm:pom-elliptic-two-time-observability-dichotomy} 的第二种情形下，
取 \(P\in E(\CC)\) 满足 \(\dd f(P)\neq 0\) 且 \(\dd(f\circ\varphi)(P)\neq 0\)。
选取局部复坐标 \(t\) 使 \(t(P)=0\)，并令
\[
F(t):=f(P(t)),\qquad G(t):=f(\varphi(P(t))).
\]
则在像曲线 \(\Phi(E)\subset \CC^2\) 的局部上存在重构映射 \(\Psi\)（\((F,G)\) 的局部逆），并满足一阶误差界
\[
|\Delta t|
\le
\frac{\bigl\|(\Delta F,\Delta G)\bigr\|}{\sqrt{|F'(t)|^2+|G'(t)|^2}}
 + o\!\left(\bigl\|(\Delta F,\Delta G)\bigr\|\right),
\]
其中 \(\|\cdot\|\) 为 \(\CC^2\) 的欧氏范数。
因此在远离临界集 \(\{\dd f=0\}\cup\{\dd(f\circ\varphi)=0\}\) 的任意紧子集上，\(\Psi\) 的局部 Lipschitz 常数由
\[
\left(\inf \sqrt{|\dd f|^2+|\dd(f\circ\varphi)|^2}\right)^{-1}
\]
上界控制，从而给出严格的误差放大率估计。
\end{proposition}
\begin{proof}
在局部坐标 \(t\) 下，\(\Phi(t)=(F(t),G(t))\) 为 \(\CC\to\CC^2\) 的 \(C^1\) 沉浸，其速度范数为
\(
|\Phi'(t)|=\sqrt{|F'(t)|^2+|G'(t)|^2}.
\)
局部逆 \(\Psi\) 定义在像曲线的邻域内，沿曲线弧长参数的一阶误差由反函数定理与均值不等式控制，
从而得到 \(|\Delta t|\le \|(\Delta F,\Delta G)\|/\inf|\Phi'|+o(\|(\Delta F,\Delta G)\|)\)。
将 \(|\Phi'|\) 以微分形式改写为 \(\sqrt{|\dd f|^2+|\dd(f\circ\varphi)|^2}\) 即得所述界。
\end{proof}
